% GENERATED FILE. Edit canonical lessons, then run npm run generate:books.
% canonical-source-hash: fnv1a64:f671b784916bc9c2
% canonical-lessons: RU-C112-problem, RU-C112-sentence, RU-C112-page, RU-C112-music, RU-C112-sport

\chapter{Problem, Sentence, Page, Music, Sport}
\label{ch:ru-problem-sentence-page-music-sport}
\hlchaptermodality{112}

\begin{chapteropening}
\textbf{By the end of this chapter:} \emph{I can say a problem, a sentence, a page, music and sport.}

Complete the last lesson of chapter 112: I can say a problem, a sentence, a page, music and sport.
\end{chapteropening}

\section[\texorpdfstring{\ru{проблема} (probléma)}{проблема (probléma)}]{\ru{проблема} (probléma) — a problem}

\label{lesson:RU-C112-problem}



\begin{quote}
Before the new one: say the Russian for homework, then the Russian for an idea.
\end{quote}

\subsection*{You'll want to know: \ru{проблема}}
\textbf{\ru{проблема}} (\emph{probléma}) — \textquotedblleft{}a problem\textquotedblright{}.

\begin{cousinweb}
\textbf{\ru{проблема}} — from Peter the Great's time; it came through Polish \emph{problema} or older German \emph{Problema} from Latin \emph{problēma}, Greek \emph{próblēma}, \textquotedblleft{}something put forward\textquotedblright{} (Vasmer).
\end{cousinweb}

\begin{grammarlens}[title={What you've built}]
One more word for this chapter.
\end{grammarlens}

\subsection*{Your turn}
\begin{itemize}
\raggedright
  \item \emph{Say it:} \emph{probléma}
  \item \emph{Say it:} \emph{probléma}, once more
  \item \emph{Say it:} say \emph{probléma}
  \item \emph{Recall:} say the Russian for homework, then the Russian for an idea, then say \emph{probléma} again
\end{itemize}

\begin{tcolorbox}[breakable,colback=teal!4,colframe=teal!35!black,title={Before you move on}]
What does \ru{проблема} mean? (\textquotedblleft{}A problem\textquotedblright{}.) Say it once more.
\end{tcolorbox}
\section[\texorpdfstring{\ru{предложение} (predlozhéniye)}{предложение (predlozhéniye)}]{\ru{предложение} (predlozhéniye) — a sentence}

\label{lesson:RU-C112-sentence}



\begin{quote}
Before the new one: say the Russian for an idea, then the Russian for a problem.
\end{quote}

\subsection*{You'll want to know: \ru{предложение}}
\textbf{\ru{предложение}} (\emph{predlozhéniye}) — \textquotedblleft{}a sentence\textquotedblright{}.

\begin{cousinweb}
\textbf{\ru{предложение}} is \emph{\ru{пред}-}, \textquotedblleft{}before\textquotedblright{}, with \emph{\ru{ложение}}, \textquotedblleft{}laying\textquotedblright{}, from \emph{-\ru{ложить}}, \textquotedblleft{}to lay\textquotedblright{}: a sentence is a thing laid before the listener; it also means \textquotedblleft{}an offer\textquotedblright{}. The root is that of \emph{\ru{лежать}}, \textquotedblleft{}to lie\textquotedblright{}, a cousin of English \emph{lie} (Vasmer).
\end{cousinweb}

\begin{grammarlens}[title={What you've built}]
One more word for this chapter.
\end{grammarlens}

\subsection*{Your turn}
\begin{itemize}
\raggedright
  \item \emph{Say it:} \emph{predlozhéniye}
  \item \emph{Say it:} \emph{predlozhéniye}, once more
  \item \emph{Say it:} say \emph{predlozhéniye}
  \item \emph{Recall:} say the Russian for an idea, then the Russian for a problem, then say \emph{predlozhéniye} again
\end{itemize}

\begin{tcolorbox}[breakable,colback=teal!4,colframe=teal!35!black,title={Before you move on}]
What does \ru{предложение} mean? (\textquotedblleft{}A sentence\textquotedblright{}.) Say it once more.
\end{tcolorbox}
\section[\texorpdfstring{\ru{страница} (stranítsa)}{страница (stranítsa)}]{\ru{страница} (stranítsa) — a page}

\label{lesson:RU-C112-page}



\begin{quote}
Before the new one: say the Russian for a problem, then the Russian for a sentence.
\end{quote}

\subsection*{You'll want to know: \ru{страница}}
\textbf{\ru{страница}} (\emph{stranítsa}) — \textquotedblleft{}a page\textquotedblright{}.

\begin{cousinweb}
\textbf{\ru{страница}} is built on \emph{\ru{страна}}, the Church Slavonic shape of \emph{\ru{сторона}}, \textquotedblleft{}side\textquotedblright{}: a page is one side of a leaf.
\end{cousinweb}

\begin{grammarlens}[title={What you've built}]
One more word for this chapter.
\end{grammarlens}

\subsection*{Your turn}
\begin{itemize}
\raggedright
  \item \emph{Say it:} \emph{stranítsa}
  \item \emph{Say it:} \emph{stranítsa}, once more
  \item \emph{Say it:} say \emph{stranítsa}
  \item \emph{Recall:} say the Russian for a problem, then the Russian for a sentence, then say \emph{stranítsa} again
\end{itemize}

\begin{tcolorbox}[breakable,colback=teal!4,colframe=teal!35!black,title={Before you move on}]
What does \ru{страница} mean? (\textquotedblleft{}A page\textquotedblright{}.) Say it once more.
\end{tcolorbox}
\section[\texorpdfstring{\ru{музыка} (múzyka)}{музыка (múzyka)}]{\ru{музыка} (múzyka) — music}

\label{lesson:RU-C112-music}



\begin{quote}
Before the new one: say the Russian for a sentence, then the Russian for a page.
\end{quote}

\subsection*{You'll want to know: \ru{музыка}}
\textbf{\ru{музыка}} (\emph{múzyka}) — \textquotedblleft{}music\textquotedblright{}.

\begin{cousinweb}
\textbf{\ru{музыка}} — Old Russian said \emph{\ru{мусикия}}; the form \emph{\ru{музыка}} came through Ukrainian and Polish \emph{muzyka} from Latin \emph{mūsica}, Greek \emph{mousik\'{\={e}}}, \textquotedblleft{}the art of the Muses\textquotedblright{} (Vasmer).
\end{cousinweb}

\begin{grammarlens}[title={What you've built}]
One more word for this chapter.
\end{grammarlens}

\subsection*{Your turn}
\begin{itemize}
\raggedright
  \item \emph{Say it:} \emph{múzyka}
  \item \emph{Say it:} \emph{múzyka}, once more
  \item \emph{Say it:} say \emph{múzyka}
  \item \emph{Recall:} say the Russian for a sentence, then the Russian for a page, then say \emph{múzyka} again
\end{itemize}

\begin{tcolorbox}[breakable,colback=teal!4,colframe=teal!35!black,title={Before you move on}]
What does \ru{музыка} mean? (\textquotedblleft{}Music\textquotedblright{}.) Say it once more.
\end{tcolorbox}
\section[\texorpdfstring{\ru{спорт} (sport)}{спорт (sport)}]{\ru{спорт} (sport) — sport}

\label{lesson:RU-C112-sport}



\begin{quote}
Before the new one: say the Russian for a page, then the Russian for music.
\end{quote}

\subsection*{You'll want to know: \ru{спорт}}
\textbf{\ru{спорт}} (\emph{sport}) — \textquotedblleft{}sport\textquotedblright{}.

\begin{cousinweb}
\textbf{\ru{спорт}} comes from English \emph{sport}, cut from \emph{disport}, \textquotedblleft{}amusement\textquotedblright{}, from Old French \emph{desporter}, \textquotedblleft{}to carry away (from work)\textquotedblright{} (Vasmer).
\end{cousinweb}

\begin{grammarlens}[title={What you've built}]
One more word for this chapter.
\end{grammarlens}

\subsection*{Your turn}
\begin{itemize}
\raggedright
  \item \emph{Say it:} \emph{sport}
  \item \emph{Say it:} \emph{sport}, once more
  \item \emph{Say it:} say \emph{sport}
  \item \emph{Recall:} say the Russian for a page, then the Russian for music, then say \emph{sport} again
\end{itemize}

\begin{tcolorbox}[breakable,colback=teal!4,colframe=teal!35!black,title={Before you move on}]
What does \ru{спорт} mean? (\textquotedblleft{}Sport\textquotedblright{}.) Say it once more.
\end{tcolorbox}
